\documentclass[aspectratio=169]{beamer}
\usepackage{kotex} % 한글 출력을 위한 패키지
\usetheme{Madrid}
\usecolortheme{default}

% Graphics 및 Code 포매팅을 위한 패키지
\usepackage{graphicx}
\usepackage{listings}
\usepackage{booktabs}
\usepackage{xcolor}

% Solidity 언어 정의 추가
\lstdefinelanguage{Solidity}{
    morekeywords={pragma, solidity, interface, contract, function, external, view, returns, address, uint256, bool, event, indexed},
    sensitive=true,
    morecomment=[l]{//},
    morecomment=[s]{/*}{*/},
    morestring=[b]",
    morestring=[b]'
}

% JSON 언어 정의 추가
\lstdefinelanguage{json}{
    morestring=[b]",
    morestring=[d]'
}

\lstset{
    basicstyle=\ttfamily\footnotesize,
    keywordstyle=\color{blue},
    stringstyle=\color{orange},
    commentstyle=\color{green!50!black},
    showstringspaces=false,
    breaklines=true
}

\title[14주차: NFT (ERC-721)]{14주차: NFT (Non-Fungible Token)}
\subtitle{Blockchain (블록체인) 기술}
\author[Prof. Kenny]{Prof. Kenny (케니 교수)}
\institute[KNU]{Kyungpook National University (경북대학교)}
\date{Spring 2026}

\begin{document}

\begin{frame}
    \titlepage
\end{frame}

\begin{frame}{Learning Objectives (학습 목표)}
    \begin{itemize}
        \item \textbf{NFT 이해:} Non-Fungible Token (대체불가능한 토큰)의 개념과 ERC-20과의 구조적 차이 이해
        \item \textbf{ERC-721 Standard (표준):} Interface (인터페이스)와 \texttt{tokenId} 기반 Ownership (소유권) 추적 메커니즘 파악
        \item \textbf{Metadata (메타데이터) \& IPFS:} \texttt{tokenURI} 구조와 Decentralized Storage (탈중앙화 저장소) 저장 이유
        \item \textbf{Hands-on (실습):} Remix IDE를 통한 Smart Contract (스마트 컨트랙트) 배포 및 Sepolia Testnet (세폴리아 테스트넷) Minting (민팅) 실습
    \end{itemize}
\end{frame}

\begin{frame}{1. NFT (Non-Fungible Token)란 무엇인가?}
    \textbf{NFT (Non-Fungible Token)}는 Blockchain (블록체인) 위에서 고유성이 보장되는 Digital Asset (디지털 자산) 소유권 증명입니다.

    \vspace{0.5cm}
    \begin{table}[]
        \centering
        \begin{tabular}{@{}lll@{}}
            \toprule
            \textbf{Feature (구분)} & \textbf{Fungible (ERC-20)} & \textbf{Non-Fungible (ERC-721)} \\ \midrule
            Interchangeability (대체 가능성) & 가능 (1 Token = 1 Token) & 불가능 (NFT \#1 $\neq$ NFT \#2) \\
            Identifier (식별자) & 없음 (수량으로 관리) & \texttt{tokenId} (고유 번호) \\
            Divisibility (분할 단위) & 소수점 가능 (Decimals) & 정수 단위만 (1개) \\
            Balance Tracking (잔액 추적) & \texttt{balances[address]} & \texttt{owners[tokenId]} \\ \bottomrule
        \end{tabular}
    \end{table}
\end{frame}

\begin{frame}[fragile]{2. ERC-721 Standard Interface (표준 인터페이스)}
    ERC-721 표준은 Ownership (소유권)과 Transfer (전송)를 관리하기 위한 필수 Function (함수)들을 정의합니다.
    
    \vspace{0.3cm}
    \begin{lstlisting}[language=Solidity, caption={Core IERC721 Functions}]
interface IERC721 {
    // Queries
    function balanceOf(address owner) external view returns (uint256);
    function ownerOf(uint256 tokenId) external view returns (address);

    // Transfers
    function transferFrom(address from, address to, uint256 tokenId) external;
    function safeTransferFrom(address from, address to, uint256 tokenId) external;

    // Approvals
    function approve(address to, uint256 tokenId) external;
    function setApprovalForAll(address operator, bool approved) external;
}
    \end{lstlisting}
\end{frame}

\begin{frame}{3. Ownership Tracking (소유권 추적): ERC-20 vs ERC-721}
    \begin{block}{ERC-20 Balance Model (잔액 모델)}
        Address (주소)가 보유한 Token (토큰)의 \textbf{Amount (수량)}에 초점을 맞춥니다.
        \begin{itemize}
            \item \texttt{balances[Alice] = 1000}
            \item \texttt{balances[Bob] = 500}
        \end{itemize}
    \end{block}

    \begin{block}{ERC-721 Ownership Model (소유권 모델)}
        고유한 Token (토큰)을 \textbf{누가 소유하고 있는지}에 초점을 맞춥니다.
        \begin{itemize}
            \item \texttt{owners[tokenId = 0] = Alice}
            \item \texttt{owners[tokenId = 1] = Bob}
            \item \texttt{owners[tokenId = 2] = Alice}
        \end{itemize}
        * \textit{Alice의 \texttt{balanceOf} 호출 결과는 2가 반환됩니다.}
    \end{block}
\end{frame}

\begin{frame}[fragile]{4. Metadata (메타데이터)와 tokenURI}
    On-chain (온체인)에 대용량 파일을 직접 저장하는 것은 Gas Fee (가스비)가 매우 높습니다. 따라서 Off-chain (오프체인) Metadata (메타데이터)를 가리키는 \texttt{tokenURI}만 저장합니다.

    \vspace{0.3cm}
    \begin{lstlisting}[language=json, caption={Standard Metadata JSON}]
{
  "name": "KNU NFT #42",
  "description": "KNU Blockchain Technology Class Commemorative NFT",
  "image": "ipfs://QmXxx.../image.png",
  "attributes": [
    { "trait_type": "Year", "value": "2026" },
    { "trait_type": "Course", "value": "Blockchain" }
  ]
}
    \end{lstlisting}
\end{frame}

\begin{frame}{5. 왜 IPFS (InterPlanetary File System)인가?}
    \begin{columns}
        \begin{column}{0.5\textwidth}
            \textbf{Centralized Server (중앙화 서버 - HTTP)}
            \begin{itemize}
                \item Location-addressed (위치 기반 주소 지정).
                \item \textcolor{red}{Risk (위험):} Server (서버)가 폐쇄되면 NFT Metadata (메타데이터)가 영구적으로 손실됩니다 (Broken Link).
            \end{itemize}
        \end{column}
        \begin{column}{0.5\textwidth}
            \textbf{Decentralized Storage (탈중앙화 저장소 - IPFS)}
            \begin{itemize}
                \item CID (Content Identifier)를 통한 Content-addressed (콘텐츠 기반 주소 지정).
                \item \textcolor{green}{Benefit (장점):} Immutable (불변성) 및 Distributed (분산). 콘텐츠가 1바이트라도 변경되면 CID가 완전히 달라집니다.
            \end{itemize}
        \end{column}
    \end{columns}
\end{frame}

\begin{frame}{6. Hands-on Session (실습 준비)}
    \begin{center}
        \Large \textbf{이제 실습을 시작해 봅시다!}
    \end{center}
    \vspace{0.5cm}
    \textbf{오늘의 Lab (실습) 사전 준비 확인:}
    \begin{enumerate}
        \item Google Chrome (구글 크롬) 브라우저를 엽니다.
        \item MetaMask (메타마스크) 로그인 후 \textbf{Sepolia Testnet (세폴리아 테스트넷)}으로 전환합니다.
        \item Sepolia ETH (세폴리아 이더) 잔액을 확인합니다 (부족한 경우 TA(조교)에게 테스트 코인을 일괄 지급받으세요).
        \item \texttt{https://remix.ethereum.org}에 접속합니다.
    \end{enumerate}
\end{frame}

\end{document}